\section{Experimental Setup}\label{sec:setup}

\subsection{Datasets and Partitioning}
\emph{CIFAR-10}~\cite{krizhevsky2009learning} contains 60{,}000 colour images of $32\times32$ pixels in ten balanced
classes. We pool its training and test sets and split them across $K=20$ clients with a label-skewed Dirichlet
partition, the standard way to control heterogeneity in FL benchmarks~\cite{li2022federated}: for every class, the
share each client receives is drawn from $\mathrm{Dir}(0.3\cdot\mathbf 1_K)$, and the draw is repeated until every
client holds at least 300 images. With $\alpha_{\mathrm{Dir}}=0.3$ most clients are dominated by a few classes
(Fig.~\ref{fig:partition}); client sizes range from 771 to 5{,}484 images (median 2{,}696). The only
pre-processing is per-channel normalisation; we use no data augmentation. The partition is fixed across all runs, so
every mechanism sees the same clients.

\emph{Adult}~\cite{becker1996adult} is the UCI census extract with 48{,}842 records; after removing records with
missing values, 45{,}222 remain. The task is to predict whether income exceeds 50K USD. Each client is one
occupation, which gives a natural, non-synthetic partition; occupations with fewer than 500 records are dropped,
leaving 12 clients with 976--6{,}020 records each (44{,}976 in total). Five numeric attributes are standardised (capital
gain and loss after a logarithm), workclass, marital status, relationship, race and sex are one-hot encoded, and
native country becomes a United-States indicator; the sampling weight, the redundant education label and the
occupation itself are dropped. Sex is the sensitive attribute for the group-fairness gaps. Adult is old and has
known limitations~\cite{ding2021retiring}; we use it because it is the only dataset in our setting with a labelled
sensitive attribute.

\emph{Splits.} Each client splits its own data at random into 70\% training, 10\% validation and 20\% test. The
validation split is used only to set hyperparameters; all reported numbers are on the test split. Keeping the split
inside each client means every client is evaluated on its own distribution, which is what personalisation targets,
and gives the membership-inference attack members (training records) and non-members (test records) from the same
distribution. Table~\ref{tab:datasets} summarises both datasets.

\begin{table}[!t]
\caption{Datasets (sizes after cleaning; client size = records per client before the 70/10/20 split).}
\label{tab:datasets}
\centering
\footnotesize
\maybetable{datasets}
\end{table}

\begin{figure}[!t]
\centering
\maybefigure{partition_cifar10}{\columnwidth}
\caption{Class composition of the 20 CIFAR-10 clients under the Dirichlet ($\alpha_{\mathrm{Dir}}=0.3$) partition.}
\label{fig:partition}
\end{figure}

\subsection{Models and Training}
For CIFAR-10 we use a small convolutional network: five $3\times3$ convolutional layers with 32, 32, 64, 64 and 128
channels, each followed by group normalisation (8 groups) and ReLU, max-pooling after the second, fourth and fifth
layers, global average pooling, and a linear head---141{,}354 parameters, of which the head has 1{,}290. Group
normalisation replaces batch normalisation because batch statistics mix examples within a batch and would break the
per-example sensitivity bound in (\ref{eq:dpsgd}); small, well-conditioned networks are also known to train best
under DP-SGD~\cite{tramer2021differentially}. For Adult we use a multilayer perceptron with hidden layers of 64 and
32 units and a linear head. FedPer keeps the final linear layer local in both models.

Training runs for $R=40$ rounds (30 for Adult). In each round every client takes $S=20$ private steps (10 for Adult)
with Poisson sampling rate $q=0.03$---about 63 records per batch for a median CIFAR-10 client---clipping norm $C=1$,
local learning rate $\eta=0.05$ (0.1 for Adult) and server learning rate 1. We set $\delta=10^{-5}$, below
$1/n_k$ for every client, and sweep $\eps\in\{0.3,0.5,1,2,4,8\}$ plus the non-private case; the common budget for
the mechanism comparison is $\eps_0=2$. Table~\ref{tab:hyper} lists all settings. For the
FedAvg schedule the calibrated noise multipliers are $\sigma\approx 10.5, 6.61, 3.57, 2.02, 1.25$ and $0.88$ for
$\eps=0.3$ to $8$; IFCA's loss reports raise them by about 6\% at $\eps_0$, and q-FFL's by about 2\%. Every
configuration is run with three seeds (0, 1, 2), which change model initialisation, sampling and noise; we report the
mean and standard deviation over seeds.

\begin{table}[!t]
\caption{Hyperparameters (identical for every mechanism unless stated).}
\label{tab:hyper}
\centering
\footnotesize
\begin{tabular}{p{4.2cm}p{3.6cm}}
\toprule
Setting & Value \\
\midrule
Clients $K$ / participation & 20 (CIFAR-10), 12 (Adult) / all, every round \\
Rounds $R$; private steps per round $S$ & 40; 20 (Adult: 30; 10) \\
Poisson sampling rate $q$; clipping norm $C$ & 0.03; 1.0 \\
Learning rate local / server & 0.05 (Adult 0.1) / 1.0 \\
$\delta$; budgets $\eps$ & $10^{-5}$; 0.3, 0.5, 1, 2, 4, 8, $\infty$ \\
Loss report clip $L_{\max}$; noise ratio $\kappa$ & 3.0; 30 \\
Fine-tuning share of the budget & 10\% (last 4 rounds' worth of steps) \\
Per-FedAvg $\alpha$ / $\beta$ / adaptation steps & 0.03 / 0.05 / 5 \\
IFCA cluster models $m$ & 3 \\
q-FFL $q$ & 0 (= FedAvg), 1, 5 \\
Dirichlet $\alpha_{\mathrm{Dir}}$; minimum client size & 0.3 (check run: 1.0); 300 \\
Seeds & 0, 1, 2 \\
\bottomrule
\end{tabular}
\end{table}

\subsection{Baselines}
Three baselines isolate the three effects at play: \emph{non-private FedAvg} gives the utility ceiling;
\emph{DP-FedAvg without personalisation} at the same budget isolates what personalisation contributes; and each
mechanism's \emph{non-private run} isolates what DP costs that mechanism. \emph{Local-only} DP training shows what a
client achieves without the federation at all.

\subsection{Implementation and Reproducibility}
All methods are implemented in PyTorch. Per-example gradients come from Opacus' gradient-sample
hooks~\cite{yousefpour2021opacus}, cross-checked against a vectorised \texttt{torch.func} implementation, and privacy
is accounted with Opacus' R\'enyi-DP accountant. For the sweeps, a single-process simulator runs all clients in turn
on one GPU: with up to 147 runs, process-per-client orchestration would multiply the run time without changing any
result. The same client code also runs in a Flower~\cite{beutel2020flower} deployment---one SuperLink and three
SuperNodes exchanging messages over HTTP on the local machine---used to demonstrate the system end to end. Experiments ran on a laptop GPU
(NVIDIA GeForce GTX 1650, 4~GB) under Windows; one DP run on CIFAR-10 takes 6.3--10.8 minutes (mean over seeds; IFCA, which trains three cluster models, is the
slowest), and the full set of 147 runs took about 14 hours. Code,
configuration files, per-run results and the exact package versions are available at
\url{https://github.com/chaaaait/t16-pfl-dp-fairness}.
